\subsection{An inline value is already a transformed reading}

Return to the source spelling \code{1,234}. At the inspected Arelle version,
an inline fact's \code{value} property first obtains assembled text and applies
the selected transformation. In the ordinary non-nil numeric path, it then
converts that result to a decimal number, applies the declared scale and
applies the sign. Both \code{textValue} and \code{stringValue} delegate to this
transformed value. Reading any of those three properties and applying scale
again would therefore change the reported amount
\citep{arelleinlinevalues2026}.

For a concrete constructed case, select \code{numdotdecimal} from the
transformation registry dated 31~July~2011, scale $-3$ and a minus sign. The
inspected function accepts \code{1,234} and removes the grouping comma, giving
1234. Scaling gives 1.234; applying the sign gives $-1.234$. If a consumer
treats that returned value as the original literal and reapplies both scale
and sign, it obtains $+0.001234$. The difference is a consumer error relative
to the selected source path. The format's namespace and function name matter:
the dispatcher selects a registry by namespace and can also invoke a custom
transformation. The example establishes no equivalence among those choices
\citep{arelleinlinevalues2026}.

Nor does \code{rawValue} recover the original source spelling in every case.
The inspected property assembles inner text, omits inline excluded content,
can follow nonnumeric continuation links, and collapses XML whitespace when a
format is present. An original spelling containing surrounding spaces and a
line break can consequently yield the same \code{1,234} reading. Escaped
content follows an additional markup-generation path with URI handling; one
inspected branch can modify an object element's \code{codebase} attribute in
the parsed tree. Original retained bytes and occurrence locators are therefore
needed even when assembled text or a later tree serialization is available
\citep{arelleinlinevalues2026}.

Consider an unescaped nonnumeric fact containing text A, an excluded child
containing X, and text B after that child, with a continuation containing C.
Under the stated traversal conditions, its assembled text is ABC. The excluded
X remains part of the source document, while B survives as the child's tail
text. The helper follows continuation links already attached by the loader;
it does not resolve the source's continuation identifier itself. If the link
is absent, a nonempty AB can still be returned. The loader separately records
missing-target, wrong-target and cycle diagnostics. Nonempty text alone cannot
establish that the disclosure was completely assembled
\citep{arelleinlinevalues2026}.

The proposed adapter thus preserves the original bytes and nodes, the ordered
text fragments and exclusions, the continuation identifiers and outcomes, and
the assembled and transformed readings. It records transform identity, scale,
sign and validation state so that each step can be reconstructed. Computing a
value is distinct from validating it: the inspected property initializes an
unvalidated state on first evaluation, and failures can set an invalid value
and raise an exception. Nil facts take a separate branch and remain distinct
from numeric zero. These source-derived requirements and constructed examples
do not replace schema checks, an executed adapter or review of an actual
issuer's facts \citep{arelleinlinevalues2026}.
